% In: settings/commands.tex

% Setup for hyperref package
\hypersetup{
    colorlinks=true,
    linkcolor=black,
    citecolor=black,
    filecolor=magenta,
    urlcolor=cyan,
}

% Path for figures (assumes a 'figures' folder in the root directory)
\graphicspath{ {./figures/} }

% Connect the .bib file to biblatex
\addbibresource{references.bib}

% Custom environment for indented quotes/citations
\newenvironment{cited3}
{\vspace{12pt}\hfill\begin{minipage}{12cm}}
        {\end{minipage}\vspace{12pt}}

% Configure the header and footer style with fancyhdr
\pagestyle{fancy}
\fancyhf{} % Clear all header and footer fields
\rhead{\fontfamily{LinuxBiolinumT-OsF}\selectfont \myheadertitle}
\lhead{\fontfamily{LinuxBiolinumT-OsF}\selectfont \footnotesize \leftmark}
\cfoot{\thepage}

\lstset{style = solarized-light, language=Python}
\renewcommand{\lstlistingname}{Código}
\renewcommand{\lstlistlistingname}{Lista de Códigos}

% Unidade imaginária
\newcommand{\iu}{\mathrm{i}}

% Números com vírgula decimal (siunitx)
\sisetup{output-decimal-marker = {,}, retain-explicit-plus = true}

% Páginas só com figuras/tabelas começam no topo, e não centralizadas
\makeatletter
\setlength{\@fptop}{0pt}
\makeatother

% Linhas de tabela geradas por codigo/gera_resultados.py (\input primitivo, que
% pode ser usado dentro de tabular sem quebrar \bottomrule)
\makeatletter
\newcommand{\linhas}[1]{\@@input tabelas/#1.tex }
\makeatother

% Trecho de um arquivo de código delimitado por comentários "# <<nome>>" e
% "# <<fim_nome>>". Uso: \trecho[opções]{arquivo.py}{nome}{legenda}
\newcommand{\trecho}[4][]{%
    \lstinputlisting[linerange={#3-fim_#3}, includerangemarker=false,
        rangeprefix=\#\ <<, rangesuffix=>>, caption={#4}, #1]{codigo/#2}}
% dracula
% black-on-white
% gnometerm
% gotham
% gruvbox-dark
% gruvbox-light
% material
% nancy
% neon
% rydgel
% nord
% seashells
% smyck
% solarized-dark
% solarized-light
% tokyonight
% white-on-black
% wild-cherry
% zenburn
% vscode
% jupyter
% mystyle
